% !TeX program = lualatex
% Standalone research notebook for original catalogue problem 67.
% Compile twice: lualatex 67.tex
% Original entry retained; research subsection and [E] references added.
% No external bibliography, image, data, or style file is required.
\documentclass[11pt,a4paper]{article}
\usepackage[margin=24mm,headheight=14pt]{geometry}
\usepackage{fontspec}
\setmainfont{Latin Modern Roman}
\setsansfont{Latin Modern Sans}
\setmonofont{Latin Modern Mono}
\usepackage{amsmath,amssymb,mathtools,amsthm}
\usepackage{unicode-math}
\setmathfont{Latin Modern Math}
\usepackage{microtype}
\usepackage{enumitem,xurl,needspace}
\usepackage[dvipsnames]{xcolor}
\usepackage{hyperref}
\hypersetup{unicode=true,colorlinks=true,linkcolor=MidnightBlue,
 urlcolor=MidnightBlue,bookmarksnumbered=true,
 pdftitle={Research notebook - Problem 67},
 pdfauthor={Research attempt; not independently refereed}}
\usepackage{bookmark}
\usepackage{titlesec}
\titleformat{\section}{\Large\bfseries\raggedright}{\thesection}{0.7em}{}
\usepackage{fancyhdr}
\pagestyle{fancy}\fancyhf{}
\fancyhead[L]{\small\sffamily Research notebook}
\fancyhead[R]{\small\sffamily Problem 67}
\fancyfoot[L]{\scriptsize 27 September 2026}
\fancyfoot[C]{\thepage}
\fancyfoot[R]{\scriptsize Unrefereed research attempt}
\setlength{\parindent}{0pt}
\setlength{\parskip}{4pt plus 1pt}
\setlength{\emergencystretch}{3em}
\clubpenalty=10000
\widowpenalty=10000
\displaywidowpenalty=10000
\setcounter{secnumdepth}{1}
\setlist[itemize]{leftmargin=3em,itemsep=2pt,topsep=3pt,parsep=0pt,before=\small}
\allowdisplaybreaks
\newcommand{\N}{\mathbb{N}}
\newcommand{\Z}{\mathbb{Z}}
\newcommand{\Q}{\mathbb{Q}}
\newcommand{\R}{\mathbb{R}}
\newcommand{\C}{\mathbb{C}}
\newcommand{\F}{\mathbb{F}}
\newcommand{\A}{\mathbb{A}}
\newcommand{\PP}{\mathbb P}
\newcommand{\E}{\mathbb{E}}
\newcommand{\eps}{\varepsilon}
\newcommand{\dd}{\,\mathrm d}
\newcommand{\sref}[2]{\hyperlink{source:#1:#2}{[S#2]}}
\newcommand{\eref}[2]{\hyperlink{extra:#1:#2}{[E#2]}}
\newif\ifshowcatalogue
\showcataloguetrue
\newcommand{\cataloguescope}[1]{\ifshowcatalogue
 \paragraph{Exact scope in the supplied catalogue.}
 \begin{quote}\small #1\end{quote}\fi}
\newtheorem{notebooklemma}{Lemma}[section]
\newtheorem{notebookproposition}[notebooklemma]{Proposition}
\numberwithin{equation}{section}
\begin{document}
\setcounter{section}{66}
% ================================================================
% problemId: problem.modularity-of-elliptic-curves-over-totally-real-fields
% source releaseRank: 67; source releaseStatus: open_with_solved_subcases
% exactTarget SHA-256: d0c44808beff43b50187e5fd77a9c816ecc7d88892fe265c23ea4e07f0f20a90
\clearpage
\section{\texorpdfstring{Modularity of Elliptic Curves over Totally Real Fields}{Modularity of Elliptic Curves over Totally Real Fields}}\label{67}
\subsection{Definitions and mathematical statement}
A number field is totally real if every embedding into $\C$ has image in $\R$. A Hilbert cuspidal newform of parallel weight two is a cuspidal Hecke eigenform whose weight is two at every real embedding, normalized and new at its level. The conjecture asks, for every elliptic curve $E/K$, for such an $f$ with
\[
 L(E,s)=L(f,s).
\]
The equality uses compatible local-factor normalizations at good and bad places, not just equality of a few computed coefficients. The field $K$ and the curve are arbitrary within the totally real scope.
\par\smallskip\textit{Formulation and concept references:} \sref{67}{1}.
\cataloguescope{For every totally real number field K and every elliptic curve E over K, prove that there is a Hilbert cuspidal newform f over K of parallel weight 2 satisfying L(E,s)=L(f,s).}
\subsection{Short English statement}
Does every elliptic curve over a totally real number field have a matching weight-two Hilbert modular form?
% BEGIN RESEARCH INSERTION
\subsection{Research attempt}

\paragraph{Current frontier.}
Freitas--Le Hung--Siksek prove modularity over real quadratic fields and finiteness of the possible nonmodular geometric isomorphism classes over each fixed totally real field \eref{67}{1}, Theorems 1 and 5. Yoshikawa's formulation source treats specified composites of real quadratic fields \sref{67}{1}; his cyclotomic-extension result retains explicit algebraic $L$-value hypotheses \eref{67}{2}, Theorem 1.1. None is replaced by a universal theorem over arbitrary totally real fields. The attempt combines finite exceptional sets with residual-image criteria and checks whether twisting or isogenies can remove the exceptional behavior.

\paragraph{Attack route.}
Potential modularity over an extension needs a descent theorem to recover the original field. A residual method needs a usable representation and a modularity theorem with the right hypotheses. A finite-exception argument needs to eliminate every exceptional geometric class, not merely almost every curve. We investigate the last route and test a natural repair: replace a difficult curve by a twist or an isogenous curve for which a small-prime modularity criterion applies.

\paragraph{Attempt: a finite elimination criterion.}
Fix $K$. The exceptional modular curves in the finite-exception proof give a finite set $J_K$ of possible $j$-invariants; no effective list or height bound is presumed \eref{67}{1}, Section 8. If $j\ne0,1728$, two $K$-curves with the same $j$ are quadratic twists. Indeed, an isomorphism over $\bar K$ has Galois discrepancy in the automorphism group $\{\pm1\}$, hence defines a quadratic character. Modularity is preserved by quadratic twisting of the associated automorphic representation, with compatible local twists at bad as well as good places. CM curves, including $j=0,1728$, are covered by the Hecke-character construction; their larger twist groups are not declared quadratic \eref{67}{1}, discussion before Lemma 11.1.

Consequently, proving modularity of one $K$-curve representing each non-CM $j\in J_K$ suffices for every elliptic curve over this fixed $K$. This is a finite conditional elimination criterion, not an algorithm computing $J_K$ or a theorem removing its remaining members. The target still requires this for every totally real $K$.

\paragraph{First escape attempt: quadratic twists.}
For a quadratic character $\chi_d$, the residual representations of $E$ and its twist satisfy
\begin{equation}\label{r67:twist}
 \bar\rho_{E^{(d)},p}(g)=\chi_d(g)\bar\rho_{E,p}(g).
\end{equation}
Thus the projective representations are the same. For any subgroup $H\le G_K$, their invariant subspaces over $\overline\F_p$ coincide: scalar multiplication by a nonzero scalar does not change whether a subspace is preserved. In particular, absolute irreducibility on $G_{K(\zeta_p)}$ is unchanged by twisting. Twists may change local reduction and conductors, so this does not say that twisting is useless; it says that this specific residual-image failure cannot be repaired in that way.

\paragraph{Second escape attempt: isogenies of arbitrary degree.}
Prime-to-$p$ isogenies identify the $p$-torsion representations directly. A $p$-power isogeny does not, but it still preserves semisimplification. Here is the needed integral-lattice argument.

\begin{notebooklemma}\label{r67:lattice}
For two $H$-stable $\Z_p$-lattices $\Lambda,\Lambda'$ in one finite-dimensional $\Q_p$-representation, the reductions modulo $p$ have isomorphic semisimplifications. This remains true after scalar extension to $\overline\F_p$.
\end{notebooklemma}
\begin{proof}
First suppose $p\Lambda\subseteq\Lambda'\subseteq\Lambda$. Set $A=\Lambda/\Lambda'$ and $B=\Lambda'/p\Lambda$. There are exact sequences
\[
 0\to B\to\Lambda/p\Lambda\to A\to0,
 \qquad
 0\to p\Lambda/p\Lambda'\to\Lambda'/p\Lambda'\to B\to0.
\]
Multiplication by $p$ identifies $A$ with $p\Lambda/p\Lambda'$. The two reductions therefore have the same composition factors with multiplicity.

For general lattices, scale one by a power of $p$ and use commensurability to arrange $p^m\Lambda\subseteq\Lambda'\subseteq\Lambda$. The successive lattices $\Lambda'+p^j\Lambda$ interpolate between them, and adjacent terms satisfy the first case. Scaling a lattice does not change its reduction up to isomorphism. Tensor the exact sequences with $\overline\F_p$ for the last assertion.
\end{proof}

A $K$-isogeny identifies rational Tate modules. Apply the lemma to their integral Tate lattices and restrict to any subgroup of $G_K$ before reducing. Absolute irreducibility on $G_{K(\zeta_p)}$ is therefore invariant under $K$-isogeny, even of degree divisible by $p$. This proves invariance of composition factors and irreducibility, not equality of full projective images or extension classes for arbitrary $p$-isogenous curves.

\paragraph{The resulting obstruction to a small-prime repair.}
The criteria in \eref{67}{1}, Theorems 3--4, give modularity when the indicated restriction is absolutely irreducible for a suitable $p\in\{3,5,7\}$. If a curve fails all three tests, every curve reached by any sequence of $K$-isogenies and quadratic twists still fails all three. Thus the proposed procedure ``change the curve until a small-prime test succeeds'' cannot remove that exceptional absolute-reducibility condition.

Failure of a sufficient criterion is not nonmodularity. For example,
\[
 E:\quad y^2+xy+y=x^3
\]
has discriminant $-26$. The tangent at $(0,0)$ is $y=0$, meeting the cubic there with multiplicity three. The chord-and-tangent law gives $3(0,0)=O$, and $(0,0)\ne O$, so the point has order three. Hence $E[3]$ has a rational invariant line, while the curve is modular over $\Q$. The same invariant-line obstruction at three persists under the operations above, without making any curve in this example nonmodular. The script checks the discriminant and the tangent substitution exactly.

\paragraph{Stress test.}
Finiteness supplies neither an effective enumeration of $J_K$ nor proofs for its remaining points. Enlarging $K$ changes the desired conclusion and introduces descent. Absolute irreducibility at a large prime by itself does not prove residual modularity over an arbitrary totally real field. Special $j$-values have larger automorphism groups and were handled separately as CM, not folded into the quadratic-twist argument. The lattice lemma preserves semisimplified residual representations but not a preferred integral lattice.

No finite set of matching Euler coefficients is used as a certificate of full modularity. A successful final elimination would need an actual comparison theorem for compatible Galois or automorphic representations and the bad-place local factors. The notebook supplies no such identification for an arbitrary residual exceptional point.

\paragraph{Outcome.}
\textbf{Outcome: Conditional reduction.}

The finite-$j$ criterion specifies what would suffice for each field, and the twisting and lattice arguments show why one natural residual-image repair cannot deliver it. The outstanding work is a different modularity argument for every exceptional geometric class, or an exhaustive elimination on the modular curves with valid stopping information. No new general modularity theorem or counterexample is claimed; the finite-exception theorem and the general lattice principle are established mathematics, explicitly applied and proved here where needed.

% END RESEARCH INSERTION
\subsection{Sources}
\begin{itemize}\raggedright
\item[\textnormal{[S1]}]\hypertarget{source:67:1}{} Sho Yoshikawa, "On the modularity of elliptic curves over a composite field of some real quadratic fields," Research in Number Theory 2 (2016), article 31.
\url{https://arxiv.org/pdf/1607.05549}.
{\small\textit{Catalogue formulation source.}}
\item[\textnormal{[S2]}]\hypertarget{source:67:2}{} Modularity of elliptic curves over cyclotomic Z\_p-extensions of real quadratic fields.
\url{https://arxiv.org/html/2206.12860v1}.
{\small\textit{Catalogue research/status reference; not a proof endorsement.}}
\item[\textnormal{[S3]}]\hypertarget{source:67:3}{} Reliability of elliptic curve data over number fields (reviewed).
\url{https://www.lmfdb.org/knowledge/show/dq.ecnf.reliability}.
{\small\textit{Catalogue research/status reference; not a proof endorsement.}}

% BEGIN ADDED RESEARCH REFERENCES
\item[\textnormal{[E1]}]\hypertarget{extra:67:1}{} Nuno Freitas, Bao V. Le Hung, and Samir Siksek, \emph{Elliptic curves over real quadratic fields are modular}, Inventiones Mathematicae 201 (2015), 159--206; arXiv:1310.7088, Theorems 1 and 3--5, Sections 8 and 11.
\url{https://arxiv.org/html/1310.7088}.
{\small\textit{Real quadratic modularity, small-prime criteria, finite geometric exceptions, and CM scope. Consulted 27 September 2026.}}
\item[\textnormal{[E2]}]\hypertarget{extra:67:2}{} Sho Yoshikawa, \emph{Modularity of elliptic curves over cyclotomic $\Z_p$-extensions of real quadratic fields}, arXiv:2206.12860 (2022), Theorem 1.1.
\url{https://arxiv.org/abs/2206.12860}.
{\small\textit{The specified algebraic $L$-value hypothesis is retained; not general totally real modularity. Consulted 27 September 2026.}}
% END ADDED RESEARCH REFERENCES
\end{itemize}

\end{document}
